\chapter{RESULTS AND DISCUSSION}

\section{General}

The seismic response of the six structural blocks of the Infectious and Communicable Disease Hospital was evaluated for fixed-base, hard-, medium- and soft-soil support conditions. Modal characteristics and response-spectrum demands were considered alongside the foundation reactions and member forces obtained from the gravity-load combination. Comparisons were made separately for each block because the structural dimensions and dynamic properties differ among the blocks.

Displacements, inter-storey drift and shear forces were evaluated for the relevant horizontal response-spectrum cases. Foundation settlements and structural member forces were assessed under $1.2DL+1.5LL$, providing complementary information on the effect of foundation flexibility under gravity loading.



\section{Modal Characteristics}

Fundamental periods and cumulative modal mass participation indicate the influence of foundation flexibility on the vibration characteristics of the six blocks. Table~\ref{tab:seq54_modal} presents the first-mode periods, while Table~\ref{tab:seq54_mass} presents the first mode achieving at least 90\% cumulative participation in both horizontal directions.

\begin{longtable}{crrrrr}
\caption{First-mode periods under the four foundation-support assumptions}\label{tab:seq54_modal}\\
\toprule
Block & Fixed (s) & Hard (s) & Medium (s) & Soft (s) & $\Delta T_{\rm soft}$ (\%)\\
\midrule
\endfirsthead
\multicolumn{6}{l}{\small Table continued}\\
\toprule
Block & Fixed (s) & Hard (s) & Medium (s) & Soft (s) & $\Delta T_{\rm soft}$ (\%)\\
\midrule
\endhead
\midrule
\multicolumn{6}{r}{\small Continued on next page}\\
\endfoot
\bottomrule
\endlastfoot
1 & 0.856 & 0.861 & 0.878 & 0.929 & 8.53 \\
2 & 0.583 & 0.594 & 0.628 & 0.698 & 19.73 \\
3 & 0.958 & 0.961 & 0.972 & 1.008 & 5.22 \\
4 & 0.723 & 0.727 & 0.743 & 0.788 & 8.99 \\
5 & 0.667 & 0.670 & 0.684 & 0.727 & 9.00 \\
6 & 0.752 & 0.756 & 0.769 & 0.811 & 7.85 \\
\end{longtable}

\begin{longtable}{crrrr}
\caption{First mode reaching at least 90\% cumulative participation in both horizontal directions}\label{tab:seq54_mass}\\
\toprule
Block & Fixed & Hard & Medium & Soft\\
\midrule
\endfirsthead
\multicolumn{5}{l}{\small Table continued}\\
\toprule
Block & Fixed & Hard & Medium & Soft\\
\midrule
\endhead
\midrule
\multicolumn{5}{r}{\small Continued on next page}\\
\endfoot
\bottomrule
\endlastfoot
1 & 9 & 9 & 9 & 8 \\
2 & 8 & 8 & 7 & 6 \\
3 & 5 & 5 & 5 & 5 \\
4 & 8 & 10 & 8 & 5 \\
5 & 6 & 6 & 6 & 6 \\
6 & 7 & 8 & 7 & 5 \\
\end{longtable}

The fundamental period increases progressively from the fixed-base to the soft-soil condition in all six blocks. The greatest percentage increase occurs in Block 2 (19.73\%), and the smallest in Block 3 (5.22\%). The number of modes required to attain 90\% cumulative mass participation ranges from 5 to 10, demonstrating differences in the dynamic characteristics of the blocks.


\begin{figure}[htbp]\centering
\includegraphics[width=0.92\textwidth,height=0.36\textheight,keepaspectratio]{figures/ch4/seq058/fundamental_period.pdf}
\caption{First-mode period for each hospital block under fixed, hard-, medium- and soft-support conditions (source: modal-analysis values in Table~\ref{tab:seq54_modal}).}\label{fig:seq55_fundamental_period}
\end{figure}
\FloatBarrier



\begin{figure}[htbp]\centering
\includegraphics[width=0.92\textwidth,height=0.36\textheight,keepaspectratio]{figures/ch4/seq058/modal_mass_90.pdf}
\caption{Earliest mode satisfying 90\% cumulative participation in both horizontal directions for each block (source: modal-analysis values in Table~\ref{tab:seq54_mass}).}\label{fig:seq55_modal_mass_90}
\end{figure}
\FloatBarrier


\section{Response-Spectrum Storey Displacement}

The uppermost-storey response-spectrum displacements were determined separately for the X and Y directions. Table~\ref{tab:seq54_disp} presents the corresponding values under the four foundation-support assumptions. As the highest modelled storey differs between blocks, the comparisons are made within individual blocks.

\begin{longtable}{clrr}
\caption{Uppermost-storey response-spectrum displacement from the six hospital blocks (mm)}\label{tab:seq54_disp}\\
\toprule
Block & Support & RS-X / X (mm) & RS-Y / Y (mm)\\
\midrule
\endfirsthead
\multicolumn{4}{l}{\small Table continued}\\
\toprule
Block & Support & RS-X / X (mm) & RS-Y / Y (mm)\\
\midrule
\endhead
\midrule
\multicolumn{4}{r}{\small Continued on next page}\\
\endfoot
\bottomrule
\endlastfoot
1 & Fixed & 59.39 & 69.63 \\
1 & Hard & 62.62 & 69.94 \\
1 & Medium & 75.37 & 70.57 \\
1 & Soft & 110.92 & 75.59 \\
2 & Fixed & 51.64 & 43.41 \\
2 & Hard & 53.24 & 46.02 \\
2 & Medium & 57.72 & 53.96 \\
2 & Soft & 67.57 & 70.97 \\
3 & Fixed & 98.88 & 73.13 \\
3 & Hard & 99.60 & 73.67 \\
3 & Medium & 101.62 & 75.12 \\
3 & Soft & 106.32 & 79.92 \\
4 & Fixed & 49.99 & 48.84 \\
4 & Hard & 50.52 & 49.44 \\
4 & Medium & 52.49 & 51.63 \\
4 & Soft & 58.50 & 57.95 \\
5 & Fixed & 66.95 & 24.09 \\
5 & Hard & 67.98 & 25.24 \\
5 & Medium & 71.60 & 28.59 \\
5 & Soft & 81.92 & 35.62 \\
6 & Fixed & 48.45 & 48.68 \\
6 & Hard & 48.93 & 49.66 \\
6 & Medium & 50.77 & 52.79 \\
6 & Soft & 56.60 & 60.07 \\
\end{longtable}

The X-direction uppermost-storey displacement increases as foundation support becomes more flexible in each of the six blocks. The increase is especially pronounced in Block 1, from 59.39~mm for the fixed-base model to 110.92~mm for soft-soil support. In Block 3, the corresponding increase is relatively small, from 98.88 to 106.32~mm. The results therefore indicate that displacement amplification depends on both foundation stiffness and the dynamic characteristics of the structure.


\begin{figure}[htbp]\centering
\includegraphics[width=0.92\textwidth,height=0.36\textheight,keepaspectratio]{figures/ch4/seq058/roof_disp_x.pdf}
\caption{Uppermost-storey displacement, X-direction response-spectrum case; comparison by block and foundation support (source: Table~\ref{tab:seq54_disp}).}\label{fig:seq55_roof_disp_x}
\end{figure}
\FloatBarrier



\begin{figure}[htbp]\centering
\includegraphics[width=0.92\textwidth,height=0.36\textheight,keepaspectratio]{figures/ch4/seq058/roof_disp_y.pdf}
\caption{Uppermost-storey displacement, Y-direction response-spectrum case; comparison by block and foundation support (source: Table~\ref{tab:seq54_disp}).}\label{fig:seq55_roof_disp_y}
\end{figure}
\FloatBarrier



\par\medskip
\begin{figure}[htbp]
\centering
\includegraphics[width=0.88\textwidth,height=0.245\textheight,keepaspectratio]{figures/appendix/B1_disp_seq057_table_derived.png}
\caption{Displacement response of Block 1.}
\label{fig:seq57_b1_displacement}
\end{figure}
\begin{figure}[htbp]
\centering
\includegraphics[width=0.88\textwidth,height=0.245\textheight,keepaspectratio]{figures/appendix/B2_disp_seq057_table_derived.png}
\caption{Displacement response of Block 2.}
\label{fig:seq57_b2_displacement}
\end{figure}
\begin{figure}[htbp]
\centering
\includegraphics[width=0.88\textwidth,height=0.245\textheight,keepaspectratio]{seq051_assets/source_plots/B3_displacement_1_8338a0.png}
\caption{Displacement response of Block 3.}
\label{fig:seq51_13}
\end{figure}
\begin{figure}[htbp]
\centering
\includegraphics[width=0.88\textwidth,height=0.245\textheight,keepaspectratio]{seq051_assets/source_plots/B4_displacement_1_497ed9.png}
\caption{Displacement response of Block 4.}
\label{fig:seq51_19}
\end{figure}
\begin{figure}[htbp]
\centering
\includegraphics[width=0.88\textwidth,height=0.245\textheight,keepaspectratio]{seq051_assets/source_plots/B5_displacement_1_87b5e8.png}
\caption{Displacement response of Block 5.}
\label{fig:seq51_25}
\end{figure}
\begin{figure}[htbp]
\centering
\includegraphics[width=0.88\textwidth,height=0.245\textheight,keepaspectratio]{seq051_assets/source_plots/B6_displacement_1_b811f4.png}
\caption{Displacement response of Block 6.}
\label{fig:seq51_31}
\end{figure}
\FloatBarrier

\section{Response-Spectrum Inter-Storey Drift}

The maximum inter-storey drift ratio provides an indication of lateral deformation demand. Table~\ref{tab:seq54_drift} compares the maximum drift ratios from the X- and Y-direction response-spectrum analyses, expressed as percentages.

\begin{longtable}{clrr}
\caption{Maximum response-spectrum inter-storey drift ratio by block and support (\%)}\label{tab:seq54_drift}\\
\toprule
Block & Support & RS-X / X (\%) & RS-Y / Y (\%)\\
\midrule
\endfirsthead
\multicolumn{4}{l}{\small Table continued}\\
\toprule
Block & Support & RS-X / X (\%) & RS-Y / Y (\%)\\
\midrule
\endhead
\midrule
\multicolumn{4}{r}{\small Continued on next page}\\
\endfoot
\bottomrule
\endlastfoot
1 & Fixed & 0.454 & 0.467 \\
1 & Hard & 0.479 & 0.462 \\
1 & Medium & 0.531 & 0.446 \\
1 & Soft & 0.662 & 0.448 \\
2 & Fixed & 0.349 & 0.279 \\
2 & Hard & 0.358 & 0.289 \\
2 & Medium & 0.383 & 0.320 \\
2 & Soft & 0.432 & 0.390 \\
3 & Fixed & 0.727 & 0.575 \\
3 & Hard & 0.730 & 0.577 \\
3 & Medium & 0.743 & 0.588 \\
3 & Soft & 0.774 & 0.620 \\
4 & Fixed & 0.530 & 0.524 \\
4 & Hard & 0.536 & 0.532 \\
4 & Medium & 0.561 & 0.562 \\
4 & Soft & 0.631 & 0.649 \\
5 & Fixed & 0.572 & 0.209 \\
5 & Hard & 0.581 & 0.216 \\
5 & Medium & 0.611 & 0.235 \\
5 & Soft & 0.695 & 0.308 \\
6 & Fixed & 0.480 & 0.455 \\
6 & Hard & 0.487 & 0.463 \\
6 & Medium & 0.513 & 0.489 \\
6 & Soft & 0.596 & 0.545 \\
\end{longtable}

The highest maximum drift ratio is 0.774\% in Block 3 under soft-soil support in the X direction. This falls within the study's Immediate Occupancy drift-screening interval, although a full performance evaluation also requires checking member demands and other damage measures. Drift changes are not identical in the two directions, reflecting differences in lateral stiffness and response.


\begin{figure}[htbp]\centering
\includegraphics[width=0.92\textwidth,height=0.36\textheight,keepaspectratio]{figures/ch4/seq058/max_drift_x.pdf}
\caption{Maximum inter-storey drift, X-direction response-spectrum case (source: Table~\ref{tab:seq54_drift}).}\label{fig:seq55_max_drift_x}
\end{figure}
\FloatBarrier



\begin{figure}[htbp]\centering
\includegraphics[width=0.92\textwidth,height=0.36\textheight,keepaspectratio]{figures/ch4/seq058/max_drift_y.pdf}
\caption{Maximum inter-storey drift, Y-direction response-spectrum case (source: Table~\ref{tab:seq54_drift}).}\label{fig:seq55_max_drift_y}
\end{figure}
\FloatBarrier



\par\medskip
\begin{figure}[htbp]
\centering
\includegraphics[width=0.88\textwidth,height=0.245\textheight,keepaspectratio]{figures/appendix/B1_drift_seq057_table_derived.png}
\caption{Drift response of Block 1.}
\label{fig:seq57_b1_drift}
\end{figure}
\begin{figure}[htbp]
\centering
\includegraphics[width=0.88\textwidth,height=0.245\textheight,keepaspectratio]{figures/appendix/B2_drift_seq057_table_derived.png}
\caption{Drift response of Block 2.}
\label{fig:seq57_b2_drift}
\end{figure}
\begin{figure}[htbp]
\centering
\includegraphics[width=0.88\textwidth,height=0.245\textheight,keepaspectratio]{seq051_assets/source_plots/B3_drift_1_7710ad.png}
\caption{Drift response of Block 3.}
\label{fig:seq51_14}
\end{figure}
\begin{figure}[htbp]
\centering
\includegraphics[width=0.88\textwidth,height=0.245\textheight,keepaspectratio]{seq051_assets/source_plots/B4_drift_1_795f60.png}
\caption{Drift response of Block 4.}
\label{fig:seq51_20}
\end{figure}
\begin{figure}[htbp]
\centering
\includegraphics[width=0.88\textwidth,height=0.245\textheight,keepaspectratio]{seq051_assets/source_plots/B5_drift_1_e52b62.png}
\caption{Drift response of Block 5.}
\label{fig:seq51_26}
\end{figure}
\begin{figure}[htbp]
\centering
\includegraphics[width=0.88\textwidth,height=0.245\textheight,keepaspectratio]{seq051_assets/source_plots/B6_drift_1_41af31.png}
\caption{Drift response of Block 6.}
\label{fig:seq51_32}
\end{figure}
\FloatBarrier

\section{Base Shear and Storey Shear}

Base shear represents the resultant horizontal reaction corresponding to each response-spectrum case. Table~\ref{tab:seq54_baseshear} presents the X- and Y-direction base shear components. The maximum storey-shear values in Table~\ref{tab:seq54_storyshear} were obtained independently from the storey-response results.

\begin{longtable}{clrr}
\caption{Response-spectrum base shear components from response-spectrum analysis (kN)}\label{tab:seq54_baseshear}\\
\toprule
Block & Support & RS-X / $|F_X|$ & RS-Y / $|F_Y|$\\
\midrule
\endfirsthead
\multicolumn{4}{l}{\small Table continued}\\
\toprule
Block & Support & RS-X / $|F_X|$ & RS-Y / $|F_Y|$\\
\midrule
\endhead
\midrule
\multicolumn{4}{r}{\small Continued on next page}\\
\endfoot
\bottomrule
\endlastfoot
1 & Fixed & 5,020.2 & 5,020.2 \\
1 & Hard & 5,020.2 & 5,020.2 \\
1 & Medium & 5,020.2 & 5,020.2 \\
1 & Soft & 5,020.2 & 5,020.3 \\
2 & Fixed & 6,578.8 & 6,578.9 \\
2 & Hard & 6,578.8 & 6,578.8 \\
2 & Medium & 6,578.8 & 6,578.9 \\
2 & Soft & 6,578.9 & 6,578.8 \\
3 & Fixed & 5,335.1 & 5,335.1 \\
3 & Hard & 5,335.2 & 5,335.1 \\
3 & Medium & 5,335.1 & 5,335.2 \\
3 & Soft & 5,335.2 & 5,335.2 \\
4 & Fixed & 1,727.5 & 1,727.5 \\
4 & Hard & 1,727.5 & 1,727.5 \\
4 & Medium & 1,727.5 & 1,727.5 \\
4 & Soft & 1,727.5 & 1,727.5 \\
5 & Fixed & 2,897.0 & 2,897.0 \\
5 & Hard & 2,897.0 & 2,897.0 \\
5 & Medium & 2,897.0 & 2,897.0 \\
5 & Soft & 2,897.0 & 2,897.0 \\
6 & Fixed & 3,465.8 & 3,465.8 \\
6 & Hard & 3,465.8 & 3,465.7 \\
6 & Medium & 3,465.8 & 3,465.7 \\
6 & Soft & 3,465.8 & 3,465.8 \\
\end{longtable}

\begin{longtable}{clrr}
\caption{Maximum response-spectrum storey-shear magnitudes from seismic analysis (kN)}\label{tab:seq54_storyshear}\\
\toprule
Block & Support & RS-X / $|V_X|$ & RS-Y / $|V_Y|$\\
\midrule
\endfirsthead
\multicolumn{4}{l}{\small Table continued}\\
\toprule
Block & Support & RS-X / $|V_X|$ & RS-Y / $|V_Y|$\\
\midrule
\endhead
\midrule
\multicolumn{4}{r}{\small Continued on next page}\\
\endfoot
\bottomrule
\endlastfoot
1 & Fixed & 5,020.2 & 5,020.2 \\
1 & Hard & 5,037.2 & 5,036.1 \\
1 & Medium & 5,049.1 & 5,048.3 \\
1 & Soft & 5,095.5 & 5,094.6 \\
2 & Fixed & 6,113.3 & 6,113.3 \\
2 & Hard & 6,133.0 & 6,198.4 \\
2 & Medium & 6,174.8 & 6,394.0 \\
2 & Soft & 6,226.3 & 6,665.9 \\
3 & Fixed & 5,335.135 & 5,335.1465 \\
3 & Hard & 5,342.7652 & 5,341.9844 \\
3 & Medium & 5,376.0873 & 5,371.8234 \\
3 & Soft & 5,412.0406 & 5,461.871 \\
4 & Fixed & 1,727.506 & 1,727.505 \\
4 & Hard & 1,732.1 & 1,732.082 \\
4 & Medium & 1,749.104 & 1,748.586 \\
4 & Soft & 1,795.269 & 1,791.967 \\
5 & Fixed & 2,897.0 & 2,897.0 \\
5 & Hard & 2,927.3 & 2,926.8 \\
5 & Medium & 3,014.0 & 2,999.0 \\
5 & Soft & 3,190.6 & 3,102.2 \\
6 & Fixed & 3,465.8 & 3,465.8 \\
6 & Hard & 3,474.3 & 3,505.5 \\
6 & Medium & 3,505.3 & 3,629.5 \\
6 & Soft & 3,588.2 & 3,855.6 \\
\end{longtable}

The storey-shear comparisons use the researcher-corrected X- and Y-direction values for Blocks 1, 3, and 4 in the same support-condition order. The storey-shear components are distinct from the base-reaction components: neither quantity is inferred from the other. For Block 3, the X-direction maximum increases from 5,335.135~kN under fixed support to 5,412.0406~kN under soft support, while the corresponding Y-direction maxima are 5,335.1465 and 5,461.871~kN. The base-shear results show little variation with foundation-support stiffness in several blocks, even though displacement and natural period increase. The storey-shear results exhibit more pronounced variation for some blocks. Consequently, foundation flexibility cannot be assumed to reduce seismic force demand in every case.


\begin{figure}[htbp]\centering
\includegraphics[width=0.92\textwidth,height=0.36\textheight,keepaspectratio]{figures/ch4/seq058/base_shear_x.pdf}
\caption{X-direction response-spectrum base shear by block and foundation condition (source: Table~\ref{tab:seq54_baseshear}).}\label{fig:seq55_base_shear_x}
\end{figure}
\FloatBarrier



\begin{figure}[htbp]\centering
\includegraphics[width=0.92\textwidth,height=0.36\textheight,keepaspectratio]{figures/ch4/seq058/story_shear_x.pdf}
\caption{Maximum X-direction response-spectrum storey shear by block and foundation condition (source: Table~\ref{tab:seq54_storyshear}).}\label{fig:seq55_story_shear_x}
\end{figure}
\FloatBarrier


\begin{figure}[htbp]\centering
\includegraphics[width=0.92\textwidth,height=0.38\textheight,keepaspectratio]{figures/ch4/seq058/story_shear_y.pdf}
\caption{Maximum Y-direction response-spectrum storey shear by block and foundation condition (source: Table~\ref{tab:seq54_storyshear}).}
\label{fig:seq057_story_shear_y}
\end{figure}
\FloatBarrier

\section{Foundation Settlement and Vertical Reaction}

Vertical foundation displacement and reaction were examined under the gravity-load combination $1.2DL+1.5LL$. Table~\ref{tab:seq54_foundation} presents the maximum absolute vertical joint displacement and the corresponding vertical base reaction for each block and support condition.

\begin{longtable}{clrr}
\caption{Maximum absolute base-joint movement and vertical reaction under 1.2DL + 1.5LL}\label{tab:seq54_foundation}\\
\toprule
Block & Support & $|U_z|$ (mm) & $F_Z$ (kN)\\
\midrule
\endfirsthead
\multicolumn{4}{l}{\small Table continued}\\
\toprule
Block & Support & $|U_z|$ (mm) & $F_Z$ (kN)\\
\midrule
\endhead
\midrule
\multicolumn{4}{r}{\small Continued on next page}\\
\endfoot
\bottomrule
\endlastfoot
1 & Fixed & 0.000 & 41,317.9 \\
1 & Hard & 0.592 & 41,317.9 \\
1 & Medium & 2.837 & 41,317.9 \\
1 & Soft & 9.925 & 41,317.9 \\
2 & Fixed & 0.552$^{a}$ & 52,026.4$^{a}$ \\
2 & Hard & 0.592 & 52,026.4 \\
2 & Medium & 2.668 & 52,026.4 \\
2 & Soft & 9.542 & 52,026.4 \\
3 & Fixed & 0.000$^{b}$ & 52,026.4$^{b}$ \\
3 & Hard & 0.556 & 39,534.4 \\
3 & Medium & 3.117 & 39,534.4 \\
3 & Soft & 11.079 & 39,534.4 \\
4 & Fixed & 0.000 & 12,916.4 \\
4 & Hard & 0.488 & 12,916.4 \\
4 & Medium & 2.360 & 12,916.4 \\
4 & Soft & 8.437 & 12,916.4 \\
5 & Fixed & 0.000 & 22,486.8 \\
5 & Hard & 0.498$^{c}$ & 0.0$^{c}$ \\
5 & Medium & 2.400 & 22,486.8 \\
5 & Soft & 8.556 & 22,486.8 \\
6 & Fixed & 0.000 & 27,242.0 \\
6 & Hard & 0.490 & 27,242.0 \\
6 & Medium & 2.381 & 27,242.0 \\
6 & Soft & 8.602 & 27,242.0 \\
\end{longtable}

Absolute vertical movement generally increases as the foundation becomes more flexible. For example, the maximum base-joint displacement in Block 1 increases from 0.592~mm under hard-soil support to 9.925~mm under soft-soil support. The values describe the analysed gravity-load condition; they are not seismic settlements. The vertical reactions show different magnitudes among the six structural blocks. Some individual reaction records are unsuitable for direct comparison and are therefore not used to infer a general reaction trend.


\begin{figure}[htbp]\centering
\includegraphics[width=0.92\textwidth,height=0.36\textheight,keepaspectratio]{figures/ch4/seq058/foundation_settlement.pdf}
\caption{Maximum foundation settlement under the $1.2DL+1.5LL$ combination (source: Table~\ref{tab:seq54_foundation}).}\label{fig:seq55_foundation_settlement}
\end{figure}
\FloatBarrier



\begin{figure}[htbp]\centering
\includegraphics[width=0.92\textwidth,height=0.36\textheight,keepaspectratio]{figures/ch4/seq058/foundation_reaction.pdf}
\caption{Vertical foundation reaction under $1.2DL+1.5LL$; provenance-flagged entries excluded (source: Table~\ref{tab:seq54_foundation}).}\label{fig:seq55_foundation_reaction}
\end{figure}
\FloatBarrier


\section{Beam and Column Force Envelopes}

The beam and column force envelopes under $1.2DL+1.5LL$ were determined by identifying the maximum absolute force or moment component among the analysed member sections. Table~\ref{tab:seq54_force} summarizes the beam shear, beam bending moment, column axial force and column bending moments. The governing member may differ between support conditions.

\begin{longtable}{clrrrrr}
\caption{Member-force absolute envelopes for the original 1.2DL + 1.5LL combination}\label{tab:seq54_force}\\
\toprule
Block & Support & Beam $|V_2|$ & Beam $|M_3|$ & Column $|P|$ & Column $|M_2|$ & Column $|M_3|$\\
\midrule
\endfirsthead
\multicolumn{7}{l}{\small Table continued}\\
\toprule
Block & Support & Beam $|V_2|$ & Beam $|M_3|$ & Column $|P|$ & Column $|M_2|$ & Column $|M_3|$\\
\midrule
\endhead
\midrule
\multicolumn{7}{r}{\small Continued on next page}\\
\endfoot
\bottomrule
\endlastfoot
1 & Fixed & 233.4 & 251.4 & 4,019.3 & 269.3 & 295.0 \\
1 & Hard & 234.4 & 247.5 & 3,998.7 & 270.2 & 299.3 \\
1 & Medium & 255.7 & 258.0 & 3,918.4 & 273.8 & 310.8 \\
1 & Soft & 426.4 & 316.7 & 3,679.8 & 283.5 & 338.3 \\
2 & Fixed & 210.5 & 230.4 & 3,739.8 & 269.7 & 266.4 \\
2 & Hard & 211.3 & 228.9 & 3,728.7 & 271.3 & 268.7 \\
2 & Medium & 214.3 & 222.5 & 3,684.9 & 277.0 & 277.5 \\
2 & Soft & 257.9 & 244.8 & 3,537.9 & 293.2 & 303.4 \\
3 & Fixed & 415.4 & 766.7 & 5,021.4 & 670.3 & 689.2 \\
3 & Hard & 415.4 & 764.9 & 5,002.7 & 671.1 & 689.8 \\
3 & Medium & 415.7 & 764.0 & 4,919.9 & 675.0 & 692.6 \\
3 & Soft & 416.9 & 774.4 & 4,694.6 & 686.8 & 700.8 \\
4 & Fixed & 298.1 & 391.9 & 2,360.9 & 140.4 & 374.1 \\
4 & Hard & 297.8 & 391.3 & 2,354.6 & 143.7 & 374.1 \\
4 & Medium & 296.8 & 388.4 & 2,328.3 & 155.5 & 374.1 \\
4 & Soft & 292.7 & 380.5 & 2,234.3 & 189.4 & 372.2 \\
5 & Fixed & 218.0 & 243.9 & 2,892.7 & 243.7 & 267.7 \\
5 & Hard & 218.1 & 241.0 & 2,881.9 & 245.8 & 269.8 \\
5 & Medium & 221.4 & 237.1 & 2,841.3 & 253.9 & 278.1 \\
5 & Soft & 231.1 & 261.7 & 2,719.0 & 280.7 & 303.4 \\
6 & Fixed & 270.5 & 274.7 & 2,843.0 & 206.7 & 238.3 \\
6 & Hard & 270.2 & 273.6 & 2,838.8 & 216.9 & 240.2 \\
6 & Medium & 267.7 & 266.7 & 2,818.4 & 249.4 & 246.8 \\
6 & Soft & 258.2 & 248.9 & 2,733.9 & 321.0 & 265.5 \\
\end{longtable}

Beam and column force envelopes vary with the foundation-support condition. Block 1, for example, shows an increase in the beam-shear envelope from 233.4~kN in the fixed-base condition to 426.4~kN under soft-soil support. The force changes are not uniform across all six blocks. These envelopes apply to the specified gravity combination and do not represent earthquake-induced force envelopes.


\begin{figure}[htbp]\centering
\includegraphics[width=0.92\textwidth,height=0.36\textheight,keepaspectratio]{figures/ch4/seq058/beam_shear_envelope.pdf}
\caption{Maximum absolute beam shear $|V_2|$ under $1.2DL+1.5LL$ (source: Table~\ref{tab:seq54_force}).}\label{fig:seq55_beam_shear_envelope}
\end{figure}
\FloatBarrier



\begin{figure}[htbp]\centering
\includegraphics[width=0.92\textwidth,height=0.36\textheight,keepaspectratio]{figures/ch4/seq058/beam_moment_envelope.pdf}
\caption{Maximum absolute beam bending moment $|M_3|$ under $1.2DL+1.5LL$ (source: Table~\ref{tab:seq54_force}).}\label{fig:seq55_beam_moment_envelope}
\end{figure}
\FloatBarrier



\begin{figure}[htbp]\centering
\includegraphics[width=0.92\textwidth,height=0.36\textheight,keepaspectratio]{figures/ch4/seq058/column_axial_envelope.pdf}
\caption{Maximum absolute column axial force $|P|$ under $1.2DL+1.5LL$ (source: Table~\ref{tab:seq54_force}).}\label{fig:seq55_column_axial_envelope}
\end{figure}
\FloatBarrier



\begin{figure}[htbp]\centering
\includegraphics[width=0.92\textwidth,height=0.36\textheight,keepaspectratio]{figures/ch4/seq058/column_moment_envelope.pdf}
\caption{Maximum absolute column moment $|M_3|$ under $1.2DL+1.5LL$ (source: Table~\ref{tab:seq54_force}).}\label{fig:seq55_column_moment_envelope}
\end{figure}
\FloatBarrier


\section{X-Direction Fragility Assessment}

\begin{figure}[htbp]\centering
\includegraphics[width=0.92\textwidth,height=0.32\textheight,keepaspectratio]{figures/ch4/seq40_part1/ground_motion_12eq_pga_selected_gorkha_component.png}
\caption{Peak ground acceleration of the twelve selected earthquake records, including the 2015 Gorkha component}
\label{fig:fragility_input_pga}
\end{figure}
\begin{figure}[htbp]\centering
\includegraphics[width=0.92\textwidth,height=0.32\textheight,keepaspectratio]{figures/ch4/part01/ground_motion_11eq_psa.png}
\caption{Five-percent-damped pseudo-spectral acceleration spectra of the eleven comparison accelerograms}
\label{fig:fragility_input_spectra}
\end{figure}
\FloatBarrier
The selected earthquake motions are described by their peak ground acceleration in Figure~\ref{fig:fragility_input_pga}. The response spectra of the eleven comparison accelerograms are shown separately in Figure~\ref{fig:fragility_input_spectra}. The Gorkha component is included in the twelve-record PGA comparison and is presented separately as an acceleration time history in the appendix.


Fragility relationships were considered for the X-direction seismic response of the six blocks under the four foundation-support conditions. The fragility curves presented in this section illustrate the probability of exceeding the damage states defined in their respective analyses. The performance categories used for inter-storey drift interpretation are given in Table~\ref{tab:legacy_isdr_performance_bands}.

The fragility curves are illustrated by shape-preserving cubic interpolation between the evaluated probability points, retaining the intensity values and exceedance probabilities at those points. The fragility relationships differ among the six blocks and foundation-support conditions. Their interpretation depends on the intensity measure and damage-state thresholds used for each curve. The adopted drift intervals provide a framework for performance interpretation but are not substituted for the thresholds used in the original fragility calculations.


\par\medskip
\noindent The block-specific X-direction fragility curves under fixed-, hard-, medium- and soft-support conditions are presented below.\par
\begin{figure}[htbp]
\centering
\includegraphics[width=0.88\textwidth,height=0.245\textheight,keepaspectratio]{figures/fragility_seq057/B1_fragility_1_236b11.png}
\caption{X-direction seismic fragility curve for Block 1 under soft support.}
\label{fig:seq51_1}
\end{figure}
\begin{figure}[htbp]
\centering
\includegraphics[width=0.88\textwidth,height=0.245\textheight,keepaspectratio]{figures/fragility_seq057/B1_fragility_2_236b11.png}
\caption{X-direction seismic fragility curve for Block 1 under medium support.}
\label{fig:seq51_2}
\end{figure}
\begin{figure}[htbp]
\centering
\includegraphics[width=0.88\textwidth,height=0.245\textheight,keepaspectratio]{figures/fragility_seq057/B1_fragility_3_236b11.png}
\caption{X-direction seismic fragility curve for Block 1 under hard support.}
\label{fig:seq51_3}
\end{figure}
\begin{figure}[htbp]
\centering
\includegraphics[width=0.88\textwidth,height=0.245\textheight,keepaspectratio]{figures/fragility_seq057/B1_fragility_4_236b11.png}
\caption{X-direction seismic fragility curve for Block 1 under fixed support.}
\label{fig:seq51_4}
\end{figure}
\begin{figure}[htbp]
\centering
\includegraphics[width=0.88\textwidth,height=0.245\textheight,keepaspectratio]{figures/fragility_seq057/B2_fragility_1_c408f3.png}
\caption{X-direction seismic fragility curve for Block 2 under soft support.}
\label{fig:seq51_5}
\end{figure}
\begin{figure}[htbp]
\centering
\includegraphics[width=0.88\textwidth,height=0.245\textheight,keepaspectratio]{figures/fragility_seq057/B2_fragility_2_c408f3.png}
\caption{X-direction seismic fragility curve for Block 2 under medium support.}
\label{fig:seq51_6}
\end{figure}
\begin{figure}[htbp]
\centering
\includegraphics[width=0.88\textwidth,height=0.245\textheight,keepaspectratio]{figures/fragility_seq057/B2_fragility_3_c408f3.png}
\caption{X-direction seismic fragility curve for Block 2 under hard support.}
\label{fig:seq51_7}
\end{figure}
\begin{figure}[htbp]
\centering
\includegraphics[width=0.88\textwidth,height=0.245\textheight,keepaspectratio]{figures/fragility_seq057/B2_fragility_4_c408f3.png}
\caption{X-direction seismic fragility curve for Block 2 under fixed support.}
\label{fig:seq51_8}
\end{figure}
\begin{figure}[htbp]
\centering
\includegraphics[width=0.88\textwidth,height=0.245\textheight,keepaspectratio]{figures/fragility_seq057/B3_fragility_1_8ae24c.png}
\caption{X-direction seismic fragility curve for Block 3 under soft support.}
\label{fig:seq51_9}
\end{figure}
\begin{figure}[htbp]
\centering
\includegraphics[width=0.88\textwidth,height=0.245\textheight,keepaspectratio]{figures/fragility_seq057/B3_fragility_2_8ae24c.png}
\caption{X-direction seismic fragility curve for Block 3 under medium support.}
\label{fig:seq51_10}
\end{figure}
\begin{figure}[htbp]
\centering
\includegraphics[width=0.88\textwidth,height=0.245\textheight,keepaspectratio]{figures/fragility_seq057/B3_fragility_3_8ae24c.png}
\caption{X-direction seismic fragility curve for Block 3 under hard support.}
\label{fig:seq51_11}
\end{figure}
\begin{figure}[htbp]
\centering
\includegraphics[width=0.88\textwidth,height=0.245\textheight,keepaspectratio]{figures/fragility_seq057/B3_fragility_4_8ae24c.png}
\caption{X-direction seismic fragility curve for Block 3 under fixed support.}
\label{fig:seq51_12}
\end{figure}
\begin{figure}[htbp]
\centering
\includegraphics[width=0.88\textwidth,height=0.245\textheight,keepaspectratio]{figures/fragility_seq057/B4_fragility_1_cb4fcf.png}
\caption{X-direction seismic fragility curve for Block 4 under soft support.}
\label{fig:seq51_15}
\end{figure}
\begin{figure}[htbp]
\centering
\includegraphics[width=0.88\textwidth,height=0.245\textheight,keepaspectratio]{figures/fragility_seq057/B4_fragility_2_cb4fcf.png}
\caption{X-direction seismic fragility curve for Block 4 under medium support.}
\label{fig:seq51_16}
\end{figure}
\begin{figure}[htbp]
\centering
\includegraphics[width=0.88\textwidth,height=0.245\textheight,keepaspectratio]{figures/fragility_seq057/B4_fragility_3_cb4fcf.png}
\caption{X-direction seismic fragility curve for Block 4 under hard support.}
\label{fig:seq51_17}
\end{figure}
\begin{figure}[htbp]
\centering
\includegraphics[width=0.88\textwidth,height=0.245\textheight,keepaspectratio]{figures/fragility_seq057/B4_fragility_4_cb4fcf.png}
\caption{X-direction seismic fragility curve for Block 4 under fixed support.}
\label{fig:seq51_18}
\end{figure}
\begin{figure}[htbp]
\centering
\includegraphics[width=0.88\textwidth,height=0.245\textheight,keepaspectratio]{figures/fragility_seq057/B5_fragility_1_18a944.png}
\caption{X-direction seismic fragility curve for Block 5 under soft support.}
\label{fig:seq51_21}
\end{figure}
\begin{figure}[htbp]
\centering
\includegraphics[width=0.88\textwidth,height=0.245\textheight,keepaspectratio]{figures/fragility_seq057/B5_fragility_2_18a944.png}
\caption{X-direction seismic fragility curve for Block 5 under medium support.}
\label{fig:seq51_22}
\end{figure}
\begin{figure}[htbp]
\centering
\includegraphics[width=0.88\textwidth,height=0.245\textheight,keepaspectratio]{figures/fragility_seq057/B5_fragility_3_18a944.png}
\caption{X-direction seismic fragility curve for Block 5 under hard support.}
\label{fig:seq51_23}
\end{figure}
\begin{figure}[htbp]
\centering
\includegraphics[width=0.88\textwidth,height=0.245\textheight,keepaspectratio]{figures/fragility_seq057/B5_fragility_4_18a944.png}
\caption{X-direction seismic fragility curve for Block 5 under fixed support.}
\label{fig:seq51_24}
\end{figure}
\begin{figure}[htbp]
\centering
\includegraphics[width=0.88\textwidth,height=0.245\textheight,keepaspectratio]{figures/fragility_seq057/B6_fragility_1_eb7bad.png}
\caption{X-direction seismic fragility curve for Block 6 under soft support.}
\label{fig:seq51_27}
\end{figure}
\begin{figure}[htbp]
\centering
\includegraphics[width=0.88\textwidth,height=0.245\textheight,keepaspectratio]{figures/fragility_seq057/B6_fragility_2_eb7bad.png}
\caption{X-direction seismic fragility curve for Block 6 under medium support.}
\label{fig:seq51_28}
\end{figure}
\begin{figure}[htbp]
\centering
\includegraphics[width=0.88\textwidth,height=0.245\textheight,keepaspectratio]{figures/fragility_seq057/B6_fragility_3_eb7bad.png}
\caption{X-direction seismic fragility curve for Block 6 under hard support.}
\label{fig:seq51_29}
\end{figure}
\begin{figure}[htbp]
\centering
\includegraphics[width=0.88\textwidth,height=0.245\textheight,keepaspectratio]{figures/fragility_seq057/B6_fragility_4_eb7bad.png}
\caption{X-direction seismic fragility curve for Block 6 under fixed support.}
\label{fig:seq51_30}
\end{figure}
\FloatBarrier

\section{Overall Discussion}

The response parameters show that foundation flexibility influences the dynamic and deformation characteristics of the hospital blocks to differing degrees. The effect on fundamental period is systematic, whereas the effects on structural forces are dependent on the response parameter and loading condition.


\subsection*{Consistency of response comparisons}
The interpretation distinguishes response-spectrum results from the separate gravity-combination outputs. Values are compared only within the corresponding block, support condition, response component and loading case; horizontal and vertical quantities are not combined. The foundation and member-force comparisons employ the $1.2DL+1.5LL$ load combination, while seismic displacements, drift and shear are evaluated using the respective response-spectrum cases. The response comparisons represent the specified structural models and support assumptions rather than direct measurements of the completed hospital.

Across all six blocks, reduced foundation restraint is associated with greater fundamental periods and higher X-direction uppermost-storey displacements. Maximum drifts also vary with soil stiffness and structural configuration. In contrast, base shear and member forces do not follow a single monotonic pattern. The comparison highlights the importance of examining displacement, drift, shear, foundation movement and member demand together when assessing soil--structure interaction. The fragility results complement these response measures by providing a probabilistic perspective on seismic performance.
